\section{Additional experiments}
\label{sec:additional_experiments}

This section provides the complete experimental results for predictive-bound
construction and population-level metric estimation, followed by ablations of
the components of \ttmethod.  We compare the following methods:
\begin{itemize}
    \item \textbf{Uncalibrated:} the raw predictive-model output, without
    conformal calibration.  In the metric-estimation experiments, this is the
    naive unweighted estimator computed from uniformly censored trajectories.
    \item \textbf{Static:} the optimized static allocation
    of~\citet{davidov2026calibrated}.  Static assigns a censoring horizon before
    observing the trajectory and does not reallocate unused turns after an
    early event.
    \item \textbf{\ttmethod:} our history-adaptive allocation, including the
    CRC budget-calibration step.
    \item \textbf{\ttmethod (w/o CRC):} the same learned adaptive policy using
    its raw projected continuation probabilities, without CRC.
    \item \textbf{Oracle:} calibration with fully observed trajectories and
    unit weights.  For metric estimation, this oracle uses the full-budget
    calibration fold; the red dashed line in the metric figures is the fixed
    benchmark value computed from the complete calibration--test population.
\end{itemize}

We repeat every experiment over 50 random calibration--test splits.  Unless
stated otherwise, the target LPB coverage is $90\%$ and the maximum trajectory
length is $\maxl=200$.  Numerical summaries in the text average over the four
target LLMs: Llama 3.1 8B Instruct, Phi-4 Mini Instruct, Qwen 2.5 14B
Instruct, and Gemma 3 4B Instruct.  We use $N_1:=|\calIone|$ throughout the paper.  The
configuration field called \texttt{n1} in the implementation instead records
the total adaptive-calibration size $|\calIone|+|\calIcrc|$.  Thus, for example,
the standard total size of 50 corresponds to $N_1=25$ and
$|\calIcrc|=25$ for \ttmethod, whereas \ttmethod without CRC uses all 50
samples in $\calIone$.  Static, Uncalibrated, and Oracle do not use this split.

For Static, ``Budget Used per Sample'' is the assigned budget
$n^{-1}\sum_i C_i$, even if an event occurs before $C_i$.  For both variants of
\ttmethod, it is the realized number of generated turns, including policy
fitting, CRC, and deployment.  Coverage variance is the sample variance across
the 50 splits in squared percentage points; bound-size variance is normalized
to one for Static separately for every dataset--target-model configuration.

\subsection{Lower predictive bounds}
\label{sec:additional_lpb_experiments}

For every LPB dataset, we report coverage, LPB value, absolute coverage
deviation, budget, number of observed events, mean inverse-probability weight,
coverage variance, and normalized LPB-value variance.  All experiments use
$N_1=25$ and $|\calIcrc|=25$ for \ttmethod.  The budget is $\bar B=20$ on
Toxicity and $\bar B=10$ on the other datasets.

\subsubsection{Toxicity}

Figures~\ref{fig:toxicity_full1}--\ref{fig:toxicity_full2} show the Toxicity
results with Detoxify as the judge.  Averaged over the four target models,
\ttmethod reduces coverage variance from $3.74$ to $1.12\,\mathrm{pp}^2$ and
reduces the mean absolute coverage deviation from $1.54$ to $0.81$ percentage
points relative to Static.  It also observes approximately 533 events per
split, compared with 181 for Static, while using 18.86 of the 20 available
turns per sample.  Its normalized LPB-value variance is $0.17$, versus one for
Static.  The no-CRC variant uses essentially the full budget and is slightly
less variable, but does not carry the CRC budget guarantee.

\begin{figure}[ht]
    \centering
    \includegraphics[width=0.9\linewidth]{figures/paper/full/lpb/toxicity_coverage_boxplot.jpg}\\
    \includegraphics[width=0.9\linewidth]{figures/paper/full/lpb/toxicity_size_boxplot.jpg}\\
    \includegraphics[width=0.9\linewidth]{figures/paper/full/lpb/toxicity_coverage_diff_boxplot.jpg}\\
    \includegraphics[width=0.9\linewidth]{figures/paper/full/lpb/toxicity_budget_used_boxplot.jpg}
    \caption{\textbf{Toxicity LPBs}: coverage, LPB value, coverage deviation,
    and budget across four target LLMs.  The target coverage is $90\%$,
    $\bar B=20$, and \ttmethod uses $N_1=25$ and
    $|\calIcrc|=25$.}
    \label{fig:toxicity_full1}
\end{figure}

\begin{figure}[ht]
    \centering
    \includegraphics[width=0.9\linewidth]{figures/paper/full/lpb/toxicity_n_observed_events_boxplot.jpg}\\
    \includegraphics[width=0.9\linewidth]{figures/paper/full/lpb/toxicity_mean_weight_boxplot.jpg}\\
    \includegraphics[width=0.9\linewidth]{figures/paper/full/lpb/toxicity_coverage_variance_barplot.jpg}\\
    \includegraphics[width=0.9\linewidth]{figures/paper/full/lpb/toxicity_size_variance_barplot.jpg}
    \caption{\textbf{Toxicity LPBs}: observed unsafe events, mean
    inverse-probability weight, coverage variance, and normalized LPB-value
    variance.  Results are computed over 50 random splits.}
    \label{fig:toxicity_full2}
\end{figure}

\subsubsection{Red Team with a Qwen judge}

Figures~\ref{fig:red_team_qwen_full1}--\ref{fig:red_team_qwen_full2} report the
Anthropic Red Team experiment evaluated by Qwen 2.5 14B Instruct.  At
$\bar B=10$, \ttmethod approximately doubles the observed unsafe events
(267 versus 137) and reduces coverage variance from $5.15$ to
$1.12\,\mathrm{pp}^2$.  Its mean coverage deviation is $0.83$ percentage
points, compared with $1.81$ for Static, and its normalized LPB-value variance
is $0.12$.  Thus, the adaptive allocation is particularly valuable in this
sparse-event regime.

\begin{figure}[ht]
    \centering
    \includegraphics[width=0.9\linewidth]{figures/paper/full/lpb/red_team_qwen_coverage_boxplot.jpg}\\
    \includegraphics[width=0.9\linewidth]{figures/paper/full/lpb/red_team_qwen_size_boxplot.jpg}\\
    \includegraphics[width=0.9\linewidth]{figures/paper/full/lpb/red_team_qwen_coverage_diff_boxplot.jpg}\\
    \includegraphics[width=0.9\linewidth]{figures/paper/full/lpb/red_team_qwen_budget_used_boxplot.jpg}
    \caption{\textbf{Red Team LPBs with a Qwen judge}: coverage, LPB value,
    coverage deviation, and budget.  The target coverage is $90\%$,
    $\bar B=10$, and \ttmethod uses $N_1=25$ and
    $|\calIcrc|=25$.}
    \label{fig:red_team_qwen_full1}
\end{figure}

\begin{figure}[ht]
    \centering
    \includegraphics[width=0.9\linewidth]{figures/paper/full/lpb/red_team_qwen_n_observed_events_boxplot.jpg}\\
    \includegraphics[width=0.9\linewidth]{figures/paper/full/lpb/red_team_qwen_mean_weight_boxplot.jpg}\\
    \includegraphics[width=0.9\linewidth]{figures/paper/full/lpb/red_team_qwen_coverage_variance_barplot.jpg}\\
    \includegraphics[width=0.9\linewidth]{figures/paper/full/lpb/red_team_qwen_size_variance_barplot.jpg}
    \caption{\textbf{Red Team LPBs with a Qwen judge}: observed unsafe events,
    mean inverse-probability weight, coverage variance, and normalized LPB-value
    variance over 50 random splits.}
    \label{fig:red_team_qwen_full2}
\end{figure}

\subsubsection{Red Team with Llama-Guard}
\label{sec:red_team_with_llamaguard}

Figures~\ref{fig:red_team_llama_guard_full1}--
\ref{fig:red_team_llama_guard_full2} use Llama-Guard as the judge.  Unsafe
events tend to occur earlier in this setup.  Consequently, all methods observe
more events than with the Qwen judge, but adaptive allocation remains useful:
\ttmethod observes approximately 1170 events, versus 590 for Static, and
reduces coverage variance from $1.63$ to $0.88\,\mathrm{pp}^2$.  Its normalized
LPB-value variance is $0.35$, and its realized budget remains below the target.

\begin{figure}[ht]
    \centering
    \includegraphics[width=0.9\linewidth]{figures/paper/full/lpb/red_team_llama_guard_coverage_boxplot.jpg}\\
    \includegraphics[width=0.9\linewidth]{figures/paper/full/lpb/red_team_llama_guard_size_boxplot.jpg}\\
    \includegraphics[width=0.9\linewidth]{figures/paper/full/lpb/red_team_llama_guard_coverage_diff_boxplot.jpg}\\
    \includegraphics[width=0.9\linewidth]{figures/paper/full/lpb/red_team_llama_guard_budget_used_boxplot.jpg}
    \caption{\textbf{Red Team LPBs with Llama-Guard}: coverage, LPB value,
    coverage deviation, and budget.  The target coverage is $90\%$,
    $\bar B=10$, and \ttmethod uses $N_1=25$ and
    $|\calIcrc|=25$.}
    \label{fig:red_team_llama_guard_full1}
\end{figure}

\begin{figure}[ht]
    \centering
    \includegraphics[width=0.9\linewidth]{figures/paper/full/lpb/red_team_llama_guard_n_observed_events_boxplot.jpg}\\
    \includegraphics[width=0.9\linewidth]{figures/paper/full/lpb/red_team_llama_guard_mean_weight_boxplot.jpg}\\
    \includegraphics[width=0.9\linewidth]{figures/paper/full/lpb/red_team_llama_guard_coverage_variance_barplot.jpg}\\
    \includegraphics[width=0.9\linewidth]{figures/paper/full/lpb/red_team_llama_guard_size_variance_barplot.jpg}
    \caption{\textbf{Red Team LPBs with Llama-Guard}: observed unsafe events,
    mean inverse-probability weight, coverage variance, and normalized LPB-value
    variance over 50 random splits.}
    \label{fig:red_team_llama_guard_full2}
\end{figure}

\subsubsection{RAG hallucinations}
\label{sec:hallucinations_exp}

In the Hallucination benchmark, the attacker constructs queries intended to
induce a response unsupported by a supplied context, and the Qwen judge assigns
the event label.  The event rate is high, so we use $\bar B=10$.  As shown in
Figures~\ref{fig:hallucination_full1}--\ref{fig:hallucination_full2}, the
advantage is smaller than in the sparse-event benchmarks but remains
consistent: \ttmethod lowers coverage variance from $0.81$ to
$0.66\,\mathrm{pp}^2$, lowers normalized LPB-value variance to $0.72$, and
observes roughly 1543 events rather than 1223.  CRC is conservative here,
using 8.17 turns per sample; the raw policy uses the full budget and observes
more events.

\begin{figure}[ht]
    \centering
    \includegraphics[width=0.9\linewidth]{figures/paper/full/lpb/hallucination_coverage_boxplot.jpg}\\
    \includegraphics[width=0.9\linewidth]{figures/paper/full/lpb/hallucination_size_boxplot.jpg}\\
    \includegraphics[width=0.9\linewidth]{figures/paper/full/lpb/hallucination_coverage_diff_boxplot.jpg}\\
    \includegraphics[width=0.9\linewidth]{figures/paper/full/lpb/hallucination_budget_used_boxplot.jpg}
    \caption{\textbf{Hallucination LPBs}: coverage, LPB value, coverage
    deviation, and budget.  The target coverage is $90\%$, $\bar B=10$, and
    \ttmethod uses $N_1=25$ and $|\calIcrc|=25$.}
    \label{fig:hallucination_full1}
\end{figure}

\begin{figure}[ht]
    \centering
    \includegraphics[width=0.9\linewidth]{figures/paper/full/lpb/hallucination_n_observed_events_boxplot.jpg}\\
    \includegraphics[width=0.9\linewidth]{figures/paper/full/lpb/hallucination_mean_weight_boxplot.jpg}\\
    \includegraphics[width=0.9\linewidth]{figures/paper/full/lpb/hallucination_coverage_variance_barplot.jpg}\\
    \includegraphics[width=0.9\linewidth]{figures/paper/full/lpb/hallucination_size_variance_barplot.jpg}
    \caption{\textbf{Hallucination LPBs}: observed hallucination events, mean
    inverse-probability weight, coverage variance, and normalized LPB-value
    variance over 50 random splits.}
    \label{fig:hallucination_full2}
\end{figure}

\subsubsection{AutoIF}
\label{sec:autoif_exp}

AutoIF evaluates whether a helper model can guide the target model to satisfy
programmatically checked instructions.  Figures~\ref{fig:autoif_full1}--
\ref{fig:autoif_full2} show the strongest LPB improvement.  At $\bar B=10$,
\ttmethod observes approximately 423 successful events, compared with 81 for
Static.  Coverage variance falls from $5.90$ to $0.68\,\mathrm{pp}^2$, mean
coverage deviation falls from $1.94$ to $0.65$ percentage points, and
normalized LPB-value variance falls to $0.03$.  The resulting performance is
close to the full-budget Oracle despite using only a small fraction of the
trajectory horizon.

\begin{figure}[ht]
    \centering
    \includegraphics[width=0.9\linewidth]{figures/paper/full/lpb/autoif_coverage_boxplot.jpg}\\
    \includegraphics[width=0.9\linewidth]{figures/paper/full/lpb/autoif_size_boxplot.jpg}\\
    \includegraphics[width=0.9\linewidth]{figures/paper/full/lpb/autoif_coverage_diff_boxplot.jpg}\\
    \includegraphics[width=0.9\linewidth]{figures/paper/full/lpb/autoif_budget_used_boxplot.jpg}
    \caption{\textbf{AutoIF LPBs}: coverage, LPB value, coverage deviation,
    and budget.  The target coverage is $90\%$, $\bar B=10$, and \ttmethod
    uses $N_1=25$ and $|\calIcrc|=25$.}
    \label{fig:autoif_full1}
\end{figure}

\begin{figure}[ht]
    \centering
    \includegraphics[width=0.9\linewidth]{figures/paper/full/lpb/autoif_n_observed_events_boxplot.jpg}\\
    \includegraphics[width=0.9\linewidth]{figures/paper/full/lpb/autoif_mean_weight_boxplot.jpg}\\
    \includegraphics[width=0.9\linewidth]{figures/paper/full/lpb/autoif_coverage_variance_barplot.jpg}\\
    \includegraphics[width=0.9\linewidth]{figures/paper/full/lpb/autoif_size_variance_barplot.jpg}
    \caption{\textbf{AutoIF LPBs}: observed successful events, mean
    inverse-probability weight, coverage variance, and normalized LPB-value
    variance over 50 random splits.}
    \label{fig:autoif_full2}
\end{figure}

\subsection{Upper predictive bounds on AutoIF}
\label{sec:autoif_upb_exp}

We also construct AutoIF UPBs at $80\%$ coverage with $\bar B=20$,
$N_1=100$, and $|\calIcrc|=100$.  Figure~\ref{fig:autoif_lpb_upb_main}
compares the LPB and UPB settings.  For UPB, \ttmethod reduces average coverage
deviation from 1.65 to 1.41 percentage points and coverage variance from 4.67
to $3.27\,\mathrm{pp}^2$, while observing 264 rather than 142 events.  The
improvement is real but smaller than for LPB: the $80\%$ UPBs lie near the
terminal horizon, leaving less room for history-adaptive decisions to change
which outcomes become identifiable.

\begin{figure}[ht]
    \centering
    \includegraphics[width=0.39\linewidth]{figures/paper/main/autoif_coverage_diff_boxplot.jpg}
    \includegraphics[width=0.39\linewidth]{figures/paper/main/autoif_budget_used_boxplot.jpg}
    \includegraphics[width=0.19\linewidth]{figures/paper/main/autoif_legend.jpg}
    \caption{\textbf{AutoIF LPB and UPB}: coverage deviation and budget.  The
    LPB target is $90\%$ with $\bar B=10$, $N_1=25$, and
    $|\calIcrc|=25$; the UPB target is $80\%$ with $\bar B=20$,
    $N_1=100$, and $|\calIcrc|=100$.}
    \label{fig:autoif_lpb_upb_main}
\end{figure}

\subsection{Population-level metric estimation}
\label{sec:unbiased_estimation}

We estimate two horizon-$\maxl$ population quantities: the event rate
$\mu_{\mathrm{event}}=\mathbb{P}(T\leq\maxl)$ and the restricted mean time to
event $\mu_{\mathrm{RMTTE}}=\mathbb{E}[\min(T,\maxl)]$.  If no event occurs by
turn $\maxl=200$, we use the sentinel event time $T=201$, so the restricted
value is 200.  The event-rate implementation uses the sequential augmented
Horvitz--Thompson (AHT) estimator.  Let $Z_i(t)=\mathbb{I}\{C_i\geq t\}$,
$\Pi_i(t)=\prod_{s=1}^{t}\rho_i(s)$, and let
$g_i(t)=\widehat{\mathbb{P}}(T_i\leq\maxl\mid H_i(t))$, with $g_i(t)$ replaced
by the resolved event indicator once the outcome is known.  The per-row AHT
contribution is
\begin{equation}
    \widehat A_i^{\mathrm{AHT}}
    =g_i(0)+\sum_{t=1}^{\maxl}
      \frac{Z_i(t)}{\Pi_i(t)}\bigl(g_i(t)-g_i(t-1)\bigr),
    \qquad
    \widehat\mu_{\mathrm{event}}
    =\frac{1}{|\calItwo|}\sum_{i\in\calItwo}
      \widehat A_i^{\mathrm{AHT}}.
    \label{eq:metric_estimator}
\end{equation}
The restricted mean uses the equivalent terminal-HT complement identity
\begin{equation}
    \widehat\mu_{\mathrm{RMTTE}}
    =\maxl-\frac{1}{|\calItwo|}\sum_{i\in\calItwo}
      \frac{\Delta_i}{\pi_i}
      \mathbb{I}\{T_i\leq\maxl\}(\maxl-T_i),
\end{equation}
where $\pi_i$ is the probability of observing the event.  Under sequential
randomization and positivity, both estimators are design-unbiased conditional
on the learned policy; correctness does not require the model probabilities
$g_i(t)$ to be calibrated.  The naive Uncalibrated estimator does not apply
this correction and can therefore be biased by right censoring.

All metric experiments use $\bar B=20$.  For \ttmethod,
$N_1=50$ and $|\calIcrc|=50$; the no-CRC variant uses all 100 adaptive-
calibration samples in $\calIone$.  Each figure reports the two estimators and
their empirical variances across splits, budget, observed events, and the
metric-target mean weighted error $A_i/\pi_i$.  The budget, event-count, and
weighted-error plots exclude Uncalibrated and Oracle because those diagnostics
are not comparable allocation quantities.  The red dashed estimator reference
is the exact value on the fully observed calibration--test population.

\subsubsection{Toxicity metrics}

Figure~\ref{fig:metrics_toxicity} shows that \ttmethod approximately halves
event-rate variance relative to Static ($7.12$ versus
$15.02\,\mathrm{pp}^2$) and reduces restricted-mean variance from 39.76 to
14.72 turns squared.  It also observes 343 events rather than 148 and lowers
the mean weighted error from 4.94 to 2.89.  The estimates remain centered near
the full-population references, whereas the naive unweighted event-rate
estimate is visibly biased downward.

\begin{figure}[ht]
    \centering
    \includegraphics[width=0.48\linewidth]{figures/paper/full/metric/dataset_toxicity/estimated_cjr_boxplot.jpg}
    \includegraphics[width=0.48\linewidth]{figures/paper/full/metric/dataset_toxicity/estimated_cjr_variance_barplot.jpg}\\
    \includegraphics[width=0.48\linewidth]{figures/paper/full/metric/dataset_toxicity/estimated_rmttu_boxplot.jpg}
    \includegraphics[width=0.48\linewidth]{figures/paper/full/metric/dataset_toxicity/estimated_rmttu_variance_barplot.jpg}\\
    \includegraphics[width=0.48\linewidth]{figures/paper/full/metric/dataset_toxicity/budget_per_sample_boxplot.jpg}
    \includegraphics[width=0.48\linewidth]{figures/paper/full/metric/dataset_toxicity/observed_jailbreaks_boxplot.jpg}\\
    \includegraphics[width=0.48\linewidth]{figures/paper/full/metric/dataset_toxicity/mean_a_weight_boxplot.jpg}
    \caption{Population-level event-rate and restricted-mean estimation on
    \textbf{Toxicity}.  Results use $\bar B=20$, $N_1=50$, and
    $|\calIcrc|=50$ and are aggregated over 50 random splits.}
    \label{fig:metrics_toxicity}
\end{figure}

\subsubsection{Red Team metrics with a Qwen judge}

On Red Team with the Qwen judge (Figure~\ref{fig:metrics_red_team_qwen}),
\ttmethod reduces event-rate variance from $17.38$ to
$10.75\,\mathrm{pp}^2$ and restricted-mean variance from 37.17 to 20.51 turns
squared.  It observes almost twice as many unsafe events as Static and lowers
the mean weighted error from 6.62 to 4.11.  The no-CRC policy is slightly less
variable in this setting and uses essentially the full target budget.

\begin{figure}[ht]
    \centering
    \includegraphics[width=0.48\linewidth]{figures/paper/full/metric/dataset_red_team_qwen/estimated_cjr_boxplot.jpg}
    \includegraphics[width=0.48\linewidth]{figures/paper/full/metric/dataset_red_team_qwen/estimated_cjr_variance_barplot.jpg}\\
    \includegraphics[width=0.48\linewidth]{figures/paper/full/metric/dataset_red_team_qwen/estimated_rmttu_boxplot.jpg}
    \includegraphics[width=0.48\linewidth]{figures/paper/full/metric/dataset_red_team_qwen/estimated_rmttu_variance_barplot.jpg}\\
    \includegraphics[width=0.48\linewidth]{figures/paper/full/metric/dataset_red_team_qwen/budget_per_sample_boxplot.jpg}
    \includegraphics[width=0.48\linewidth]{figures/paper/full/metric/dataset_red_team_qwen/observed_jailbreaks_boxplot.jpg}\\
    \includegraphics[width=0.48\linewidth]{figures/paper/full/metric/dataset_red_team_qwen/mean_a_weight_boxplot.jpg}
    \caption{Population-level metric estimation on \textbf{Red Team with a
    Qwen judge}.  The configuration is $\bar B=20$, $N_1=50$, and
    $|\calIcrc|=50$.}
    \label{fig:metrics_red_team_qwen}
\end{figure}

\subsubsection{Red Team metrics with Llama-Guard}

Figure~\ref{fig:metrics_red_team_llama_guard} exhibits the largest metric-
estimation gain.  \ttmethod lowers event-rate variance from $25.07$ to
$6.85\,\mathrm{pp}^2$ and restricted-mean variance from 73.79 to 13.17 turns
squared.  It observes approximately 1216 events, compared with 252 for Static,
and lowers the mean weighted error from 8.49 to 2.73.  CRC is conservative in
realized budget (15.59 turns per sample), while the no-CRC policy remains near
20 and is correspondingly somewhat less variable.

\begin{figure}[ht]
    \centering
    \includegraphics[width=0.48\linewidth]{figures/paper/full/metric/dataset_red_team_llama_guard/estimated_cjr_boxplot.jpg}
    \includegraphics[width=0.48\linewidth]{figures/paper/full/metric/dataset_red_team_llama_guard/estimated_cjr_variance_barplot.jpg}\\
    \includegraphics[width=0.48\linewidth]{figures/paper/full/metric/dataset_red_team_llama_guard/estimated_rmttu_boxplot.jpg}
    \includegraphics[width=0.48\linewidth]{figures/paper/full/metric/dataset_red_team_llama_guard/estimated_rmttu_variance_barplot.jpg}\\
    \includegraphics[width=0.48\linewidth]{figures/paper/full/metric/dataset_red_team_llama_guard/budget_per_sample_boxplot.jpg}
    \includegraphics[width=0.48\linewidth]{figures/paper/full/metric/dataset_red_team_llama_guard/observed_jailbreaks_boxplot.jpg}\\
    \includegraphics[width=0.48\linewidth]{figures/paper/full/metric/dataset_red_team_llama_guard/mean_a_weight_boxplot.jpg}
    \caption{Population-level metric estimation on \textbf{Red Team with
    Llama-Guard}.  The configuration is $\bar B=20$, $N_1=50$, and
    $|\calIcrc|=50$.}
    \label{fig:metrics_red_team_llama_guard}
\end{figure}

\subsubsection{Hallucination metrics}

On Hallucination (Figure~\ref{fig:metrics_hallucination}), \ttmethod reduces
event-rate variance from $27.80$ to $10.11\,\mathrm{pp}^2$ and restricted-mean
variance from 91.82 to 10.39 turns squared.  CRC is notably conservative here,
using 11.62 turns per sample.  The raw policy uses 19.94 turns and approaches
the full-budget Oracle variance (event-rate variance $0.88\,\mathrm{pp}^2$),
showing both the quality of the learned allocation and the finite-sample cost
of the rigorous budget correction in this high-event-rate regime.

\begin{figure}[ht]
    \centering
    \includegraphics[width=0.48\linewidth]{figures/paper/full/metric/dataset_hallucination/estimated_cjr_boxplot.jpg}
    \includegraphics[width=0.48\linewidth]{figures/paper/full/metric/dataset_hallucination/estimated_cjr_variance_barplot.jpg}\\
    \includegraphics[width=0.48\linewidth]{figures/paper/full/metric/dataset_hallucination/estimated_rmttu_boxplot.jpg}
    \includegraphics[width=0.48\linewidth]{figures/paper/full/metric/dataset_hallucination/estimated_rmttu_variance_barplot.jpg}\\
    \includegraphics[width=0.48\linewidth]{figures/paper/full/metric/dataset_hallucination/budget_per_sample_boxplot.jpg}
    \includegraphics[width=0.48\linewidth]{figures/paper/full/metric/dataset_hallucination/observed_jailbreaks_boxplot.jpg}\\
    \includegraphics[width=0.48\linewidth]{figures/paper/full/metric/dataset_hallucination/mean_a_weight_boxplot.jpg}
    \caption{Population-level metric estimation on the
    \textbf{Hallucination} benchmark.  The configuration is $\bar B=20$,
    $N_1=50$, and $|\calIcrc|=50$.}
    \label{fig:metrics_hallucination}
\end{figure}

\subsubsection{AutoIF metrics}

Figure~\ref{fig:metrics_autoif} shows that the improvement extends to a utility
rather than safety event.  \ttmethod reduces event-rate variance from $16.05$
to $8.18\,\mathrm{pp}^2$ and restricted-mean variance from 45.70 to 17.69 turns
squared.  It more than doubles the observed successful events and lowers the
mean weighted error from 4.73 to 2.71.  Thus, the allocation principle is not
specific to unsafe-event detection.

\begin{figure}[ht]
    \centering
    \includegraphics[width=0.48\linewidth]{figures/paper/full/metric/dataset_autoif/estimated_cjr_boxplot.jpg}
    \includegraphics[width=0.48\linewidth]{figures/paper/full/metric/dataset_autoif/estimated_cjr_variance_barplot.jpg}\\
    \includegraphics[width=0.48\linewidth]{figures/paper/full/metric/dataset_autoif/estimated_rmttu_boxplot.jpg}
    \includegraphics[width=0.48\linewidth]{figures/paper/full/metric/dataset_autoif/estimated_rmttu_variance_barplot.jpg}\\
    \includegraphics[width=0.48\linewidth]{figures/paper/full/metric/dataset_autoif/budget_per_sample_boxplot.jpg}
    \includegraphics[width=0.48\linewidth]{figures/paper/full/metric/dataset_autoif/observed_jailbreaks_boxplot.jpg}\\
    \includegraphics[width=0.48\linewidth]{figures/paper/full/metric/dataset_autoif/mean_a_weight_boxplot.jpg}
    \caption{Population-level event-rate and restricted-mean estimation on
    \textbf{AutoIF}.  The configuration is $\bar B=20$, $N_1=50$, and
    $|\calIcrc|=50$.}
    \label{fig:metrics_autoif}
\end{figure}

Across all five benchmarks, \ttmethod with CRC has lower empirical variance
than Static for both the event-rate and restricted-mean estimators.  The
no-CRC version is often slightly more efficient because it avoids the
finite-sample CRC correction, but only the CRC variant is designed to provide
the distribution-free expected-budget guarantee.  These results also clarify
why the full-budget Oracle cannot generally be matched: it resolves essentially
every trajectory and therefore removes acquisition noise, whereas any method
with $\bar B\ll\maxl$ must retain some inverse-probability variability.

\subsection{Ablation study for \ttmethod}
\label{sec:ablation_exps}

All ablations use the Toxicity dataset with Qwen 2.5 14B Instruct as both
attacker and target and Detoxify as the judge, except for the attacker-shift
study.  The default cell has $\bar B=20$ and total adaptive-calibration size
$|\calIone|+|\calIcrc|=50$; hence \ttmethod uses $N_1=25$ and
$|\calIcrc|=25$.  Every plot includes Static, \ttmethod, and \ttmethod without
CRC.  For LPB ablations, the weighted-error diagnostic is always the selected-
target quantity $A_i(q_{\widehat\tau})/\pi_i$; we do not select among several
post-hoc definitions of $A_i$.

Following the task-selection rule used for the paper, we show metric-estimation
results for the optimization-process, score-noise, score-function, and
representation-capacity ablations.  The calibration-size, budget, optimization-
objective, CRC-cap, and attacker-shift ablations use the LPB task.

\subsubsection{Adaptive-calibration sample size (LPB)}

Figure~\ref{fig:ablation_n1} varies the total adaptive-calibration size over
$50,100,200,300,$ and $400$.  For \ttmethod with CRC these correspond to
$(N_1,|\calIcrc|)=(25,25),(50,50),(100,100),(150,150),$ and $(200,200)$;
without CRC, the entire displayed total is $N_1$.  Performance is stable from
50 through 200 total samples.  At 300--400 samples, the larger fitting cost
leaves fewer deployment resources, the number of observed events falls, and
both weighted error and coverage variance rise.  Thus, more Phase-I data is
not automatically better; the default total of 50 gives the best overall
exploration--deployment tradeoff in this experiment.

\begin{figure}[ht]
    \centering
    \includegraphics[width=0.98\linewidth]{figures/paper/ablations/dapro_n1_ablation.jpg}
    \caption{LPB adaptive-calibration-size ablation.  The x-axis is
    $|\calIone|+|\calIcrc|$, not $N_1$ alone.  Lines show means and shaded
    regions show one standard deviation across 50 random splits.}
    \label{fig:ablation_n1}
\end{figure}

\subsubsection{Available budget (LPB)}

Figure~\ref{fig:ablation_budget} varies $\bar B$ from 5 to 50 and pairs each
budget with the total adaptive-calibration sizes
$25,50,50,100,150,$ and $200$, respectively.  For CRC runs, the total is split
as evenly as possible between $\calIone$ and $\calIcrc$; for example, the
$\bar B=5$ cell uses $(N_1,|\calIcrc|)=(13,12)$.  \ttmethod has markedly lower
coverage deviation, coverage variance, and selected-target weighted error than
Static at every budget.  Increasing budget reveals more events and reduces
weight variability, with diminishing returns beyond approximately 20 turns.
Both adaptive variants remain near their target budgets; CRC is mildly
conservative, especially in the smallest-budget cell.

\begin{figure}[ht]
    \centering
    \includegraphics[width=0.98\linewidth]{figures/paper/ablations/dapro_budget_ablation.jpg}
    \caption{LPB budget ablation.  Parenthesized x-axis values give the total
    adaptive-calibration size $|\calIone|+|\calIcrc|$.}
    \label{fig:ablation_budget}
\end{figure}

\subsubsection{Optimization objective (LPB)}

Figure~\ref{fig:ablation_objective} compares the soft-prefix objective used by
\ttmethod with hard-terminal and soft-terminal objectives.  The hard-prefix
case is omitted because its realized objective is algebraically equivalent to
the hard-terminal case.  Soft prefix is the most reliable choice: it observes
the most events and attains the lowest weighted error and coverage variance.
Hard terminal is competitive but less efficient, while soft terminal performs
substantially worse because a single terminal probability gives a weaker
learning signal for the sequence of continuation decisions.  This supports
using the prefix decomposition even though the final variance target is
terminal.

\begin{figure}[ht]
    \centering
    \includegraphics[width=0.98\linewidth]{figures/paper/ablations/dapro_hard_soft_ablation.jpg}
    \caption{LPB optimization-objective ablation at $\bar B=20$,
    $N_1=25$, and $|\calIcrc|=25$.}
    \label{fig:ablation_objective}
\end{figure}

\subsubsection{Optimization process (metric estimation)}

To isolate the effect of Phase-I optimization, Figure~
\ref{fig:ablation_optimization} compares \ttmethod with a Uniform+CRC policy
whose continuation probabilities are not optimized.  Uniform+CRC devotes all
50 adaptive-calibration samples to CRC and uses no policy-fitting fold.  It
meets the budget and observes many events, but its event-rate variance is
$47.76\,\mathrm{pp}^2$ and its mean weighted error is 15.94.  In comparison,
\ttmethod with CRC attains $10.49\,\mathrm{pp}^2$ and weighted error 2.81;
Static attains $12.00\,\mathrm{pp}^2$.  The large gap demonstrates that budget
calibration alone is insufficient: the learned allocation must control the
reach probabilities of the target contributions.  It also shows that the raw
number of observed events is not, by itself, a variance objective.

\begin{figure}[ht]
    \centering
    \includegraphics[width=0.98\linewidth]{figures/paper/ablations/dapro_metric_optimization_ablation.jpg}
    \caption{Metric-estimation optimization-process ablation at
    $\bar B=20$.  Standard \ttmethod uses $N_1=25$ and
    $|\calIcrc|=25$; Uniform+CRC uses all 50 rows for CRC.}
    \label{fig:ablation_optimization}
\end{figure}

\subsubsection{Score corruption (metric estimation)}

We corrupt the learned score according to
\begin{equation}
    S_i^{(\lambda)}(t)
    =(1-\lambda)\,\mathcal{S}(H_i(t))+\lambda U_i(t),
    \qquad U_i(t)\sim\operatorname{Unif}(0,1).
\end{equation}
Thus, larger $\lambda$ means more noise and $\lambda=1$ is fully random.
Figure~\ref{fig:ablation_score_noise} shows that the CRC policy is stable for
$\lambda\leq0.75$, but at $\lambda=1$ its event-rate variance increases to
$15.23\,\mathrm{pp}^2$, the event count decreases, and the weighted error
increases.  The no-CRC curve is nonmonotone across finite splits, although its
event count declines and weighted error rises at high noise.  The main
conclusion is therefore not that the score is irrelevant: moderate corruption
is absorbed by the coarse $K=2$ projection and CRC rescaling, whereas removing
all signal measurably harms the guaranteed policy.

\begin{figure}[ht]
    \centering
    \includegraphics[width=0.98\linewidth]{figures/paper/ablations/dapro_metric_score_noise_ablation.jpg}
    \caption{Metric-estimation score-noise ablation at $\bar B=20$,
    $N_1=25$, and $|\calIcrc|=25$.}
    \label{fig:ablation_score_noise}
\end{figure}

\subsubsection{Score function (metric estimation)}

Figure~\ref{fig:ablation_score} compares the estimated current event
probability, the estimated remaining-time quantile, a random score, and the
oracle remaining time while holding $K=2$ and the budget controller fixed.
For \ttmethod with CRC, the estimated remaining-time score gives the lowest
event-rate error and a slightly lower variance than the current-event score;
the current-event score nevertheless has the lowest weighted error and reveals
the most events.  Random scoring has the largest CRC variance.  Oracle
remaining time does not dominate because remaining-time ordering is not the
oracle ordering for the event-rate variance objective.  Since rankings also
change slightly without CRC, no alternative score uniformly dominates; the
current-event probability is the most robust default across diagnostics.

\begin{figure}[ht]
    \centering
    \includegraphics[width=0.98\linewidth]{figures/paper/ablations/dapro_metric_score_ablation.jpg}
    \caption{Metric-estimation score-function ablation at $\bar B=20$,
    $N_1=25$, $|\calIcrc|=25$, and $K=2$.}
    \label{fig:ablation_score}
\end{figure}

\subsubsection{Representation capacity (metric estimation)}

Figures~\ref{fig:ablation_representation}--
\ref{fig:ablation_representation_diagnostics} compare $K\in\{1,2,4,8\}$ score
bins with a continuous monotone score map.  Increasing capacity from $K=1$
improves the fitted Phase-I and held-out Phase-II objectives, confirming that
the score contains information discarded by a constant policy.  With CRC,
$K=8$ has the lowest event-rate variance ($6.62\,\mathrm{pp}^2$), while the
continuous and $K=4$ maps are close.  Without CRC, $K=2$ is the most stable
finite-sample choice.  Hence, richer maps can improve the CRC variant, but the
gain is not monotone or universal; $K=2$ remains a defensible low-variance,
low-overfitting default shared by the main experiments.

\begin{figure}[ht]
    \centering
    \includegraphics[width=0.98\linewidth]{figures/paper/ablations/dapro_metric_representation_ablation.jpg}
    \caption{Metric-estimation representation-capacity ablation at
    $\bar B=20$, $N_1=25$, and $|\calIcrc|=25$.}
    \label{fig:ablation_representation}
\end{figure}

\begin{figure}[ht]
    \centering
    \includegraphics[width=0.98\linewidth]{figures/paper/ablations/dapro_metric_representation_diagnostics.jpg}
    \caption{Phase-I objective, held-out Phase-II weighted-error proxy, and
    event-rate variance for the representation-capacity ablation.}
    \label{fig:ablation_representation_diagnostics}
\end{figure}

\subsubsection{CRC row-cost cap (LPB)}

Figure~\ref{fig:ablation_cmax} varies $C_{\max}$ from $0.1\bar B$ to
$10\bar B=\maxl$.  The no-CRC policy is unchanged, as expected, because this
cap belongs to the CRC controller.  A very small cap prevents the policy from
following informative long trajectories and sharply underuses the budget.  A
very large cap makes the finite-sample CRC correction loose and therefore also
induces conservatism.  Values near $\bar B$--$2\bar B$ give the best tradeoff;
we use $C_{\max}=2\bar B$ because it retains low variance while allowing an
individual trajectory to receive more than the average budget.

\begin{figure}[ht]
    \centering
    \includegraphics[width=0.98\linewidth]{figures/paper/ablations/dapro_cmax_ablation.jpg}
    \caption{LPB CRC row-cost-cap ablation at $\bar B=20$,
    $N_1=25$, and $|\calIcrc|=25$.}
    \label{fig:ablation_cmax}
\end{figure}

\subsection{Robustness to attacker shift}
\label{sec:attacker_shift}

Finally, we learn and calibrate the policy on Qwen-attacker trajectories and
evaluate it on Gemma-attacker trajectories, using the Llama 3.1 target model.
We test both Red Team and Toxicity at $\bar B=10$, $N_1=25$, and
$|\calIcrc|=25$.  Toxicity keeps Detoxify fixed; in the available Red Team
collections, the LLM judge follows the attacker family, so that experiment
changes both attacker and judge.  This is a genuine distribution shift, and
exchangeability between calibration and test no longer holds; accordingly, we
do not claim the nominal-coverage theorem for this experiment.

Figure~\ref{fig:attacker_shift} shows conservative coverage under shift
($94$--$97\%$ rather than $90\%$).  Although this makes absolute coverage
deviation an unsuitable measure of adaptation quality here, \ttmethod sharply
reduces coverage variance: from $2.17$ to $0.42\,\mathrm{pp}^2$ on Red Team and
from $2.96$ to $0.13\,\mathrm{pp}^2$ on Toxicity.  It also observes substantially
more events and lowers the selected-target weighted error while remaining at
or below the target budget.  Thus, the learned allocation remains statistically
stable under this shift, even though the shifted coverage distribution is
conservative.

\begin{figure}[ht]
    \centering
    \includegraphics[width=0.98\linewidth]{figures/paper/ablations/dapro_attacker_shift_ablation.jpg}
    \caption{LPB attacker-shift ablation.  The policy is calibrated on Qwen
    attacker data and evaluated on Gemma attacker data for Red Team and
    Toxicity.  The target coverage is $90\%$ and $\bar B=10$.}
    \label{fig:attacker_shift}
\end{figure}
